\documentclass[11pt]{article}
\usepackage[letterpaper,margin=0.82in]{geometry}
\usepackage[T1]{fontenc}
\usepackage{lmodern}
\usepackage{microtype}
\usepackage{amsmath,amssymb}
\setlength{\parindent}{0pt}
\setlength{\parskip}{0.55em}
\pagestyle{empty}

\begin{document}

\begin{center}
{\LARGE\bfseries Nash's Problem}\par
{\large eight curl-free fields, and why there are exactly eight}\par
\end{center}

\textbf{problem.} find $X\subset\mathbb R^3$ such that, on $U=\mathbb R^3\setminus X$, the space $V$ of continuously differentiable vector fields with $\nabla\times F=0$, modulo the space $W$ of gradient fields $F=\nabla f$, has $\dim(V/W)=8$.

\textbf{construction.} remove eight parallel lines:
\[
L_j=\{(j,0,z):z\in\mathbb R\},\qquad
X=\bigcup_{j=1}^{8}L_j.
\]
for each $j$, define on $U$
\[
F_j(x,y,z)=\frac{1}{2\pi}
\left(
-\frac{y}{(x-j)^2+y^2},
\frac{x-j}{(x-j)^2+y^2},
0
\right).
\]
each $F_j$ is smooth on $U$, independent of $z$, and a direct differentiation gives $\nabla\times F_j=0$. locally, $F_j$ is $\nabla\theta_j/(2\pi)$, where $\theta_j$ is the angle around $L_j$; globally, that angle changes by $2\pi$ after one turn.

\textbf{eight independent classes.} let $\gamma_i(t)=(i+\varepsilon\cos t,\varepsilon\sin t,0)$ for $0\le t\le2\pi$ and $0<\varepsilon<1/3$. a direct parametrization gives $F_i(\gamma_i(t))\cdot\gamma_i'(t)=1/(2\pi)$. if $j\ne i$, the disk bounded by $\gamma_i$ misses $L_j$, so stokes' theorem gives
\[
\oint_{\gamma_i}F_j\cdot d\mathbf r=\delta_{ij}.
\]
every gradient has zero integral around a closed curve. hence no nonzero linear combination of the $F_j$ is a gradient: their eight classes in $V/W$ are independent.

\textbf{why there are no others.} sliding points vertically to $z=0$ retracts $U$ onto the plane with the eight points $(j,0)$ removed. every closed curve in this punctured plane is, for purposes of integrating a curl-free field, a sum of loops around those eight punctures. equivalently, its first homology group is $\mathbb Z^8$, generated by the $\gamma_j$.

now take any $F\in V$ and set
\[
a_j=\oint_{\gamma_j}F\cdot d\mathbf r,
\qquad
G=F-\sum_{j=1}^{8}a_jF_j.
\]
the displayed circulation identity makes every period of $G$ around $\gamma_j$ zero. since these loops generate all periods, $\oint_\gamma G\cdot d\mathbf r=0$ for every closed curve $\gamma$ in $U$. therefore $f(p)=\int_{p_0}^{p}G\cdot d\mathbf r$ is path independent and $G=\nabla f\in W$. thus
\[
[F]=\sum_{j=1}^{8}a_j[F_j]
\qquad\text{and}\qquad
\boxed{\dim(V/W)=8}.
\]

\end{document}
